% Generated draft: requires mathematical and source review.
\documentclass[11pt]{article}
\usepackage[margin=25mm]{geometry}
\usepackage{fontspec}
\setmainfont{Latin Modern Roman}
\usepackage{amsmath,amssymb,mathtools,unicode-math}
\setmathfont{Latin Modern Math}
\usepackage{enumitem,needspace,xurl,hyperref}
\setlength{\parindent}{0pt}
\setlength{\parskip}{5pt}
\setlength{\emergencystretch}{3em}
\newcommand{\N}{\mathbb{N}}
\newcommand{\Z}{\mathbb{Z}}
\newcommand{\Q}{\mathbb{Q}}
\newcommand{\R}{\mathbb{R}}
\newcommand{\C}{\mathbb{C}}
\newcommand{\F}{\mathbb{F}}
\newcommand{\A}{\mathbb{A}}
\newcommand{\PP}{\mathbb{P}}
\newcommand{\E}{\mathbb{E}}
\newcommand{\eps}{\varepsilon}
\newcommand{\dd}{\,\mathrm{d}}
\begin{document}
\section*{UnsolvedMath ID: 62}
\textbf{Source problem code:} catalogue entry
\subsection{Definitions and mathematical statement}
The supplied statement is attributed to David Hilbert (1900). Write the general seventh-degree equation, after making it monic when that normalization is permitted, as
\[
 P_a(x)=x^7+a_6x^6+a_5x^5+a_4x^4+a_3x^3+a_2x^2+a_1x+a_0=0.
\]
Here $a=(a_0,\ldots,a_6)$ denotes the variable coefficients; their domain, the domain of $x$, and the permitted parameter regions are not specified in the record. A local single-valued root branch on a parameter region $U$ would be a function $r$ satisfying $P_a(r(a))=0$ for every $a\in U$. Alternatively a solution could mean a multivalued algebraic function or an enumeration of all roots; the supplied statement does not choose between these targets.

The stated request is to prove that the general equation of degree seven cannot be solved using functions of at most two variables. To turn this into a precise impossibility assertion, one must fix the allowed function class $\mathcal F$, its domains and regularity or algebraicity requirements, the permitted compositions and other operations, and the meaning of solving the equation. Relative to such choices, the intended assertion would say that the required general root function or functions have no representation by the permitted finite compositions of members of $\mathcal F$ having at most two arguments. These choices are unresolved in the supplied record, so this is an incomplete formulation, not a uniquely specified conjecture.

The supplied literature triage, dated 2026-08-17, distinguishes the continuous-function variant, reported as answered affirmatively by the Kolmogorov--Arnold representation theorem, from an algebraic-function variant reported as unresolved. Thus the unqualified impossibility statement must not be asserted for arbitrary continuous constituent functions. The triage explicitly requires restoration of the function class and allowed compositions, and does not supply the missing algebraic formulation. The registry label is \emph{open}, whereas the dated triage label is \emph{partially solved}; these are supplied classifications, not a new determination of current mathematical status.
\subsection{Short English statement}
Hilbert asked whether the general equation of degree seven can be solved by composing functions of two variables. The record requests an impossibility proof but leaves the permitted functions and the meaning of a general solution unspecified. Its dated review says the continuous version has an affirmative answer and distinguishes an unresolved algebraic version, whose precise restrictions still need to be supplied.
\subsection{Sources}
\begin{itemize}
\item UnsolvedMath local registry, \texttt{data/unsolvedmath/problems.json}, permanent ID 62, source code \texttt{catalogue entry}, ``Hilbert's 13th Problem: Seventh Degree Equations.'' The supplied statement and background attribute the problem to David Hilbert (1900).
\item Encyclopedia of Mathematics, ``Hilbert problems,'' \url{https://encyclopediaofmath.org/wiki/Hilbert_problems}. The supplied literature triage lists this authoritative secondary source as accessed 2026-08-17 and says its status overview distinguishes the historical problem's unresolved aspects. No independent source consultation is claimed here.
\item Ulam AI and the UnsolvedMath Contributors, \emph{UnsolvedMath: A Curated Collection of Open Mathematics Problems}; local source record and its embedded literature triage dated 2026-08-17. Dataset curation and original metadata are supplied under CC BY 4.0; referenced source material retains its own terms.
\end{itemize}
\end{document}
